% GDZ PPN236018647,native540lower–541;HKUST541–542.
\NumberedItem{55}
\indent Supplementum eximium ad polygonorum descriptionem inveni. Sc[ilicet], si \(a,b,c,d,\ldotp\,\ldotp\,\ldotp\) sint factores primi numeri primi \(p\) unitate truncati, tunc ad polygoni \(p\) laterum descriptionem nihil aliud requiri quam ut:
\par\indent \(1^0.\) arcus indefinitus in \(a,b,c,d,\ldotp\,\ldotp\,\ldotp\) partes secetur,
\par\indent \(2^0.\) ut polygona \(a,b,c,d,\ldotp\,\ldotp\,\ldotp\) laterum describantur.
\par\noindent\hfill Gotting[ae, 1797] Ian. 19.

\medskip\begingroup\small
\indent Hier ist jedenfalls Folgendes gemeint. Bedeuten wie im Texte \(a,b,c,\ldotp\,\ldotp\,\ldotp\) die Primteiler von \(p-1\)
\par\endgroup
\clearpage\SourceLeaf{541}\ArithmeticHead{512}{L}{TAGEBUCH MIT ERLÄUTERUNGEN.}
\begingroup\small
\noindent und setzt man
\[
p-1=aA,\quad A=bB,\quad B=cC,\quad\ldotp\,\ldotp\,\ldotp,
\]
\noindent so bestimmen sich die \(A\)-gliedrigen Perioden nach der im art. 352. der \textit{Disquisitiones arithmeticae}, Werke I, S. 431 entwickelten Methode durch eine Gleichung \(a\)-ten Grades, die durch Auflösung von \(x^a=1\) und Teilung eines gewissen Winkels in \(a\) gleiche Teile gelöst wird; darauf die \(B\)-gliedrigen Perioden durch eine Gleichung \(b\)-ten Grades, die durch Auflösung von \(x^b=1\) und Teilung eines gewissen Winkels in \(b\) gleiche Teile gelöst wird u.s.w.; vergl. die Bemerkung zu der Nr. 66.
\par\noindent\hfill\textsc{Bachmann.}
\par\endgroup

\NumberedItem{56}
\indent Theoremata de res[iduis] \(-1,\mp2\) simili methodo demonstrata ut cetera. . . .
\par\noindent\hfill Gott[ingae, 1797] Febr. 4.

\medskip\begingroup\small
\indent Dieselbe Zeitangabe Werke I, S. 476, Anmerkung zum art. 145. der \textit{Disquisitiones arithmeticae}, Werke I, S. 110.
\par\noindent\hfill\textsc{Klein.}
\par\endgroup

\NumberedItem{57}
\indent Forma
\[
aa+bb+cc-bc-ac-ab,
\]
\noindent quod ad divisores attinet, convenit cum hac:
\[
aa+3bb.
\]
\noindent\hfill[1797] Febr. 6.

\medskip\begingroup\small
\indent Setzt man
\[
a^2+b^2+c^2-bc-ac-ab=\alpha,
\]
\noindent so ist
\[
(2a-b-c)^2+3(b-c)^2=4\alpha.
\]
\noindent Jeder ungerade Primteiler der ersten Form ist also auch Teiler der Form \(x^2+3y^2\). Umgekehrt, ist \(p\) ein Teiler der letzteren, so gibt es eine ungerade Zahl \(A\), für die
\[
A^2\equiv-3\ (\operatorname{mod.}p)
\]
\noindent ist; setzt man \(B=1\), so ist also:
\[
A^2+3B^2\equiv0\ (\operatorname{mod.}p).
\]
\noindent Wählt man dann
\[
a=0,\quad b=\frac{A+B}{2},\quad c=\frac{A-B}{2},
\]
\noindent so ist \(p\) auch Teiler der ersten Form.
\par\noindent\hfill\textsc{Bachmann.}
\par\endgroup
